% Chapter 5: Conclusion and Future Work
\section{Conclusions}
This thesis has explored the application of Distributed Source Coding (DSC) techniques to per-coordinate gradient compression (quantization) in federated learning. We have demonstrated that DSC can significantly reduce the communication overhead associated with transmitting gradients from workers to a central server, while maintaining high model accuracy and convergence speed. Our experiments examined multiple protocols, finding that simply having more side information, in the form of reconstructed gradients of other workers, is not enough, and a protocol which only has access to a single worker's gradient can outperform it. This was because the transmission cost of sending the encoder/decoder pairs was cheaper than expected, and allowed for the use of a protocol that directly trains on the gradient that will be sent. On the other hand, the protocol which needs other workers' reconstructed gradients couldn't take advantage of this, as its options for training were limited to the past gradients of other workers. A hybrid approach that combines elements of both protocols could potentially yield even better results. Even with this limitation, the All-Out protocol still achieved a respectable compression rate and maintained good accuracy and convergence speed.
